\documentclass[11pt]{article}

\usepackage[margin=1in]{geometry}
\usepackage{amsmath,amssymb,amsthm,bm}
\usepackage{booktabs}
\usepackage{hyperref}

\newtheorem{proposition}{Proposition}
\newtheorem{remark}{Remark}
\newcommand{\dd}{\mathrm{d}}
\newcommand{\ii}{\mathrm{i}}
\newcommand{\norm}[1]{\left\lVert #1 \right\rVert_2}

\title{Improving NPBS and Extending the Residual-Metric Idea to Radiative Transfer}
\author{Repository audit note}
\date{\today}

\begin{document}
\maketitle

\begin{abstract}
NPBS is best understood as a residual-correction framework tied to a chosen
left preconditioner.  The original Helmholtz construction uses the
constant-coefficient damped Green operator from the Convergent Born Series
(CBS), but the same logic applies to any computable map \(B_h\) that turns a
physical residual \(r=f-A_hu\) into a better-scaled residual coordinate
\(B_hr\).  This note gives the self-contained algebra, lists concrete
Helmholtz improvements, and develops a radiative-transfer-equation (RTE)
extension using both a transport resolvent and a transport-plus-diffusion
synthetic acceleration preconditioner.
\end{abstract}

\section{Purpose}

This note records concrete mathematical and algorithmic improvements suggested by
the repository audit of the Neural Preconditioned Born Series (NPBS) framework.
It also proposes and implements a first radiative-transfer-equation (RTE)
generalization. Here RTE means the steady-state monoenergetic radiative transfer
equation discretized by discrete ordinates.  All operators below are discrete
operators; continuum notation is used only to explain their origin.

\section{Sharper Mathematical Framing}

The main idea should be stated as a preconditioner-induced residual geometry
principle, not only as a CBS-specific construction. Let
\[
  A_h u = f
\]
be a discrete linear PDE system and let \(B_h\) be a computationally cheap
left-preconditioning map. The preconditioned residual is
\[
  z = B_h(f-A_h u),
\]
and the natural residual seminorm induced by \(B_h\) is
\[
  \|r\|_{B_h} := \|B_h r\|_2 .
\]
CBS/FFT Green preconditioning corresponds to a particular \(B_h\); classical
shifted-Laplacian multigrid, diffusion synthetic acceleration, and transport
synthetic acceleration are other possible choices.
If \(B_h\) is injective on the residual space, this seminorm is a norm.
Otherwise it is still a valid coordinate objective, but convergence in
\(\|B_hr\|_2\) must be supplemented by a lower bound relating \(B_hr\) to the
physical residual \(r\).

Given a correction model \(P_\theta\), the metric-matched one-step training
objective associated with \(B_h\) is
\[
  \mathcal{L}_{B_h}(\theta)
  =
  \mathbb{E}_{r}
  \frac{\|B_h(A_hP_\theta(B_hr)-r)\|_2}{\|B_hr\|_2}.
\]
The learned solver is then
\[
  u_{m+1}=u_m+P_\theta B_h(f-A_hu_m).
\]
Thus the choice of \(B_h\) fixes both the inference coordinates and the residual
geometry used in training.

\subsection{Discrete CBS Identity}

For Helmholtz, the paper identity should be stated at the level of the chosen
discrete split. Let
\[
  A_h = L_{\eta,h} - V_{\eta,h}, \qquad G_{\eta,h} = L_{\eta,h}^{-1}.
\]
In the repository scaling, \(A_h\) includes a factor of \(dx^2\), so the
Born-coordinate preconditioner is
\[
  B_{\eta,h} r := G_{\eta,h}\left(\frac{r}{dx^2}\right).
\]
Then
\[
  B_{\eta,h} A_h
  =
  I - G_{\eta,h} V_{\eta,h}.
\]
This is the exact algebraic statement behind the implementation. It should not
be conflated with every shifted-Laplacian preconditioner. In particular:
\[
  B_G r = G_{\eta,h}(r/dx^2)
\]
is the pure Born/constant-reference shifted Helmholtz map, while implemented CBS
uses
\[
  B_{\mathrm{CBS}} r = \Gamma_{\eta,h} B_G r,
  \qquad
  \Gamma_{\eta,h} = \ii V_{\eta,h}/\epsilon .
\]
Classical variable-coefficient shifted-Laplacian V-cycles use a different
operator \(M_\beta^{-1}\).

\subsection{Residual-Propagator Theorem}

The most useful theorem for NPBS is the residual-propagator identity.

\begin{proposition}[NPBS residual propagation]
Let \(B_h\) be a fixed linear left preconditioning map and define
\[
  z_m := B_h(f-A_h u_m).
\]
Consider the learned iteration
\[
  u_{m+1} = u_m + P_\theta z_m .
\]
Then
\[
  z_{m+1}
  =
  (I - B_h A_h P_\theta) z_m .
\]
Consequently, if the residual-coordinate map
\[
  F_\theta(z):=z-B_hA_hP_\theta(z)
\]
satisfies
\[
  \|F_\theta(z)\|_2\le \rho\|z\|_2,\qquad \rho<1,
\]
on a trajectory-invariant residual-coordinate set, then
\[
  \|z_m\|_2 \le \rho^m \|z_0\|_2 .
\]
If \(P_\theta\) is linear, this condition reduces to
\[
  \|I-B_hA_hP_\theta\|_{\mathcal{Z}}\le \rho .
\]
\end{proposition}

\begin{proof}
By the update rule,
\[
  f-A_h u_{m+1}
  =
  f-A_h u_m - A_h P_\theta z_m .
\]
Applying \(B_h\) and using \(z_m=B_h(f-A_hu_m)\) gives
\[
  z_{m+1}
  =
  z_m - B_h A_h P_\theta z_m
  =
  (I-B_hA_hP_\theta)z_m.
\]
\end{proof}

This theorem explains why the training target should not merely approximate
\(A_h^{-1}\), but should make \(B_h A_h P_\theta\) close to the identity on the
solver's residual trajectory.

\section{Immediate Improvements for Helmholtz NPBS}

\begin{enumerate}
\item \textbf{Separate operators in all tables.}
Report \(B_G\), \(B_{\mathrm{CBS}}\), and variable-coefficient shifted
Laplacian as separate algorithms.

\item \textbf{Formalize the damping policy.}
Use
\[
  k_0(a)=\omega/\operatorname{mean}(c_a), \qquad
  \eta(a)=\alpha\|\omega^2/c_a^2-k_0(a)^2\|_\infty,
\]
with \(\alpha\) selected by validation and fixed consistently between training
and inference.

\item \textbf{Train on solver-state residuals.}
Replace purely iid Gaussian probes by a mixture of Gaussian probes, point-source
probes, band-limited probes, and replay residuals collected after
\(1,2,4,8\) CBS/NPBS/FGMRES steps.

\item \textbf{Make FGMRES a first-class target.}
Report matrix-vector calls, preconditioner calls, restart length, convergence
rate, and wall time. Include CBS inner solves and both \(B_G\) and
\(B_{\mathrm{CBS}}\).

\item \textbf{Audit complex linearity.}
For a learned map \(P_\theta\), measure
\[
  \frac{\|P_\theta(\ii r)-\ii P_\theta(r)\|_2}{\|P_\theta(r)\|_2}.
\]
If this defect is large, either use FGMRES only in flexible mode or constrain
the architecture.

\item \textbf{Separate discretization claims.}
Pseudo-spectral CBS/NPBS, finite-difference shifted Laplacian, sparse shifted
LU, and direct finite-difference LU should be separate comparison regimes.
\end{enumerate}

\section{Helmholtz Implementation Audit}

The Helmholtz code uses the scaled discrete system
\[
  A_hu=b,\qquad
  A_hu = dx^2\left(G_{\eta,h}^{-1}u - V_{\eta,h}u\right).
\]
Consequently the Born-coordinate residual used in training and inference is
\[
  B_G r = G_{\eta,h}(r/dx^2).
\]
This is exactly what the repository implements in the CBS and learned-NPBS
paths.  In the file \texttt{models/cbs.py}, the stationary CBS update is
\[
  u^{m+1}
  =
  u^m-\Gamma_{\eta,h}\left[
    u^m-G_{\eta,h}\left(V_{\eta,h}u^m+b/dx^2\right)
  \right],
  \qquad
  \Gamma_{\eta,h}=\ii V_{\eta,h}/\epsilon.
\]
The implemented residual operator is
\[
  A_hu = dx^2(G_{\eta,h}^{-1}u-V_{\eta,h}u),
\]
so the identity
\[
  G_{\eta,h}\left(V_{\eta,h}u+b/dx^2\right)-u
  =
  G_{\eta,h}\left((b-A_hu)/dx^2\right)
\]
is algebraically correct.  The local verifier
\[
  \texttt{helmholtz/verify\_operator\_identities.py}
\]
checks this identity and, in double precision, gives a relative error of order
\(10^{-15}\) on the tested case.

The learned NPBS models in \texttt{models/model2d.py} use three related
training modes:
\[
\begin{array}{ccl}
\text{loss type 1} &:& \text{input }r,\quad
  \|r-A_hP_\theta(r)\|_2/\|r\|_2,\\[2mm]
\text{loss type 2} &:& \text{input }B_Gr,\quad
  \|r-A_hP_\theta(B_Gr)\|_2/\|r\|_2,\\[2mm]
\text{loss type 3} &:& \text{input }B_Gr,\quad
  \|B_G(r-A_hP_\theta(B_Gr))\|_2/\|B_Gr\|_2.
\end{array}
\]
Loss type 3 is the metric-matched NPBS objective.  The Krylov inference path
keeps the same coordinate convention: for loss types 2 and 3 it applies
\(G_{\eta,h}(r/dx^2)\) before calling the learned preconditioner.  This is the
correct right-preconditioning map for flexible GMRES.

The implementation is therefore algebraically consistent, but it is not
perfect.  The main limitations are:
\begin{enumerate}
\item The Euclidean residual norms are not quadrature weighted.  On a fixed
grid this is harmless for relative comparisons, but a continuum-consistent
norm should include spatial weights and, when appropriate, source scaling.
\item Point sources are inserted as grid entries.  A physical Dirac
approximation would require dimension-dependent scaling, for example
\(1/dx^2\) in two dimensions before the repository's \(dx^2\)-scaled operator
is formed.
\item The CBS damping parameter \(\epsilon\) has no explicit positive floor.
With the default absorbing boundary this is usually nonzero, but constant
media with zero boundary width can make \(\epsilon=0\).
\item The two-dimensional CBS class does not initialize a zero right-hand side
when setup is called without a source and then followed by a stationary solve.
The training code bypasses this by calling the learned preconditioner directly,
but the solver API should still define \(\texttt{b}\) robustly.
\item Constructors call \(\texttt{torch.set\_default\_dtype}\), which is a
global side effect.  A cleaner implementation would keep dtype local to tensors.
\item The absorbing-boundary formulas are implementation approximations, and
the ABC/PBL code paths in \texttt{models/cbs.py} and
\texttt{models/model2d.py} should be harmonized.  The verified identity is for
the operator actually assembled by the code, not for an independent physical
boundary discretization.
\item Learned updates are not forced to be contractive.  Some architectures are
not exactly complex-linear in the residual input, so flexible GMRES is the
appropriate Krylov wrapper when these maps are used as preconditioners.
\item Training with iid Gaussian right-hand sides tests an operator-learning
problem, but it does not fully match residuals produced along CBS, NPBS, or
FGMRES trajectories.  Replay residuals should be included for a stricter solver
benchmark.
\end{enumerate}

Thus the Helmholtz NPBS implementation is a coherent and useful prototype of
the paper's algebra.  It should not be described as a perfect production solver
without the robustness and benchmarking improvements above.

\section{RTE Generalization}

\subsection{Discrete-Ordinates RTE}

Consider the steady-state RTE
\[
  \Omega\cdot\nabla_x \psi(x,\Omega)
  + \sigma_t(x)\psi(x,\Omega)
  - \sigma_s(x)\int_{\mathbb{S}^{d-1}}p(\Omega,\Omega')\psi(x,\Omega')
    \,\dd \Omega'
  =
  q(x,\Omega).
\]
For discrete ordinates \(\{\Omega_\ell,w_\ell\}_{\ell=1}^{N_\Omega}\) and
isotropic scattering,
\[
  (A\psi)_\ell
  =
  \Omega_\ell\cdot\nabla \psi_\ell
  + \sigma_t \psi_\ell
  - \sigma_s \phi,
  \qquad
  \phi = \sum_m w_m\psi_m .
\]

\subsection{Transport Resolvent Coordinates}

Choose a constant reference attenuation \(\sigma_0+\eta>0\) and define
\[
  (L_{\eta}\psi)_\ell
  =
  \Omega_\ell\cdot\nabla \psi_\ell
  +(\sigma_0+\eta)\psi_\ell .
\]
On a periodic spatial grid, \(L_\eta\) is diagonal in Fourier space:
\[
  \widehat{(G_\eta r)}_\ell(\xi)
  =
  \frac{\hat r_\ell(\xi)}
  {\ii \Omega_\ell\cdot\xi + \sigma_0+\eta}.
\]
Define
\[
  V_\eta := L_\eta - A .
\]
For isotropic scattering,
\[
  (V_\eta\psi)_\ell
  =
  (\sigma_0+\eta-\sigma_t)\psi_\ell
  + \sigma_s \phi .
\]
Then the same identity holds:
\[
  I-G_\eta V_\eta = G_\eta A.
\]
Therefore the RTE has the same NPBS structure:
\[
  \psi^{m+1}
  =
  \psi^m + P_\theta\!\left(G_\eta(q-A\psi^m)\right),
\]
with the metric-matched loss
\[
  \mathcal{L}^{R_\eta}_{\mathrm{RTE}}(\theta)
  =
  \mathbb{E}_{r}
  \frac{
    \|G_\eta(A P_\theta(G_\eta r)-r)\|_2
  }{
    \|G_\eta r\|_2
  }.
\]

\subsection{General Synthetic Acceleration as Preconditioner}

The pure transport inverse \(G_\eta\) damps streaming-dominated modes but can be
weak in highly scattering regimes, where the slow error is mostly a scalar-flux
mode.  Classical diffusion synthetic acceleration (DSA) addresses this by
adding a low-order diffusion correction.

Write
\[
  T\psi := \Omega\cdot\nabla\psi+\sigma_t\psi,\qquad
  S\psi := \sigma_s\phi,\qquad
  A=T-S.
\]
Let \(B_T\) denote a transport preconditioner.  It may be the periodic
constant-reference inverse \(G_\eta\), one upwind transport sweep, several
transport sweeps, or a sparse/block transport preconditioner.  Given an angular
residual \(r\), first compute a transport correction
\[
  \delta\psi_T = B_T r.
\]
The residual left after that correction is
\[
  e_T = r-A\delta\psi_T.
\]
Synthetic acceleration constructs an additional correction from a reduced
model of this defect.  Let
\[
  R_{\rm lo}: \hbox{angular residuals}\to \hbox{low-order variables}
\]
be a restriction, let \(A_{\rm lo}\) be a low-order operator, and let
\[
  P_{\rm lo}: \hbox{low-order variables}\to \hbox{angular fluxes}
\]
be a prolongation or lift.  The one-shot synthetic-acceleration preconditioner
is
\[
  B_{\rm SA} r
  =
  B_T r
  +
  P_{\rm lo}A_{\rm lo}^{-1}R_{\rm lo}(I-AB_T)r .
\]
This formula is the most useful way to define the preconditioner \(B_h\) for
NPBS:
\[
  B_h = B_{\rm SA}.
\]
It is a residual-to-correction map.  It does not need to be an exact Green
operator and usually does not satisfy an exact Born identity
\(I-GV=GA\).

\subsection{DSA as a Special Case}

For isotropic scattering, the slow part of \(e_T\) is well represented by its
angular moment
\[
  m_T = \sum_\ell w_\ell (e_T)_\ell.
\]
The constant-coefficient diffusion-augmentation approximation solves
\[
  \left(-D_0\Delta+\sigma_{a,0}\right)\delta\phi = m_T,
  \qquad
  D_0=\frac{1}{2\sigma_{t,0}},\qquad
  \sigma_{a,0}=\sigma_{t,0}-\sigma_{s,0},
\]
where the factor \(2\) is the spatial dimension of the periodic prototype.  On
a periodic grid, this diffusion equation is diagonal in Fourier variables:
\[
  \widehat{\delta\phi}(\xi)
  =
  \frac{\widehat{m_T}(\xi)}
  {D_0|\xi|^2+\sigma_{a,0}}.
\]
Let \(\mathcal{I}\delta\phi\) be the isotropic lift satisfying
\[
  \sum_\ell w_\ell(\mathcal{I}\delta\phi)_\ell=\delta\phi .
\]
The DSA-style transport-plus-diffusion preconditioner is then
\[
  B_{\mathrm{TD}}r
  :=
  \delta\psi_T+\mathcal{I}\delta\phi .
\]
Equivalently,
\[
  B_{\mathrm{TD}}
  =
  B_T
  +
  \mathcal{I}
  \left(-D_0\Delta+\sigma_{a,0}\right)^{-1}
  \Phi(I-AB_T),
\]
where \(\Phi r=\sum_\ell w_\ell r_\ell\).  In the repository prototype,
\[
  B_T=G_\eta,\qquad
  D_0 = 1/(2\sigma_{t,0}),\qquad
  \sigma_{a,0}=\sigma_{t,0}-\sigma_{s,0},
\]
and all operators use periodic pseudo-spectral derivatives.  This is a
constant-coefficient periodic DSA-style prototype, not a full transport
discretization with boundary-consistent DSA.

\subsection{TSA-Style Preconditioner}

Transport synthetic acceleration (TSA) replaces the diffusion low-order solve
by a transport-like low-order solve.  It is useful when the diffusion
approximation is too crude, for example in optically thin regions, anisotropic
media, strongly heterogeneous media, or problems where angular structure
matters.

The abstract TSA construction keeps the same two-stage residual correction:
\[
  \delta\psi_T = B_T r,\qquad
  e_T=(I-AB_T)r.
\]
Instead of taking only the scalar flux moment, choose a transport restriction
\[
  R_{\rm TSA}e_T = e_{\rm lo},
\]
where \(e_{\rm lo}\) may contain a coarser angular quadrature, selected angular
moments, a lower-order \(P_N\) or \(S_N\) representation, or a coarser spatial
transport grid.  Let \(A_{\rm TSA}\) be the corresponding low-order transport
operator and let \(P_{\rm TSA}\) lift the low-order correction back to the
full angular grid.  A TSA-style preconditioner is
\[
  B_{\rm TSA}r
  =
  B_T r
  +
  P_{\rm TSA}A_{\rm TSA}^{-1}R_{\rm TSA}(I-AB_T)r .
\]
Examples of \(A_{\rm TSA}\) include:
\[
\begin{array}{ll}
\text{coarse angular transport:}
&
  \Omega_j^{\rm lo}\cdot\nabla \chi_j
  +\sigma_t\chi_j
  -\sigma_s\sum_k w_k^{\rm lo}\chi_k,
\\[1mm]
\text{coarse-grid transport:}
&
  A_{\rm TSA}=R_x A_{\rm transport} P_x
  \quad\text{on a coarser spatial mesh},
\\[1mm]
\text{moment transport:}
&
  A_{\rm TSA}\text{ is a closed low-order moment or quasidiffusion system}.
\end{array}
\]
The inverse \(A_{\rm TSA}^{-1}\) may itself be approximate.  With \(\nu\)
inner low-order iterations,
\[
  c_0=0,\qquad
  c_{j+1}=c_j+\widetilde B_{\rm lo}
  \left(R_{\rm TSA}e_T-A_{\rm TSA}c_j\right),
  \qquad j=0,\ldots,\nu-1,
\]
and
\[
  B_{\rm TSA}^{(\nu)}r = B_T r + P_{\rm TSA}c_\nu .
\]
This is the precise meaning of ``TSA-style \(B_h\)'' in the NPBS framework:
\(B_h\) is the full composite synthetic-acceleration map
\[
  r
  \mapsto
  B_T r
  +
  P_{\rm TSA}A_{\rm TSA}^{-1}R_{\rm TSA}(r-AB_T r),
\]
not merely the raw transport Green operator \(G_\eta\).

\begin{remark}
Unlike the pure transport resolvent \(G_\eta\), \(B_{\mathrm{TD}}\) and
\(B_{\rm TSA}\) are not presented as exact Born splits \(I-GV=GA\).  They are
synthetic-acceleration preconditioners.  The residual-propagator theorem still
applies with \(B_h=B_{\mathrm{TD}}\) or \(B_h=B_{\rm TSA}\), and the
metric-matched objective becomes
\[
  \mathcal{L}_{B_{\mathrm{TD}}}(\theta)
  =
  \mathbb{E}_{r}
  \frac{
    \|B_{\mathrm{TD}}(A P_\theta(B_{\mathrm{TD}}r)-r)\|_2
  }{
    \|B_{\mathrm{TD}}r\|_2
  }.
\]
This is the appropriate way to generalize NPBS beyond exact CBS-style Green
coordinates.  The implementation should be read as a periodic
constant-coefficient diffusion synthetic preconditioner prototype, not as a
boundary-compatible production DSA or TSA discretization.
\end{remark}

\subsection{Why This Is Promising}

The RTE already has a classical source-iteration/synthetic-acceleration
ecosystem. Classical source iteration is a Neumann iteration driven by a
variable-coefficient transport sweep and the scattering perturbation; DSA and
other synthetic accelerations are preconditioned residual corrections. NPBS can
generalize this by replacing a hand-designed low-order correction with a
learned correction trained in the metric induced by the chosen transport or
synthetic-acceleration preconditioner.

The most promising RTE version is not a direct end-to-end intensity predictor.
It is a learned residual corrector used inside source iteration or FGMRES:
\[
  B_h \in \{G_{\eta,\mathrm{transport}},\; B_{\mathrm{DSA}},\;
  B_{\mathrm{SA}},\; \hbox{hybrids}\}.
\]
This also suggests a useful RTE experiment: compare source iteration, DSA,
FGMRES+DSA, NPBS with transport \(G_\eta\), and NPBS with DSA-induced metric on
scattering-dominated media.

\begin{remark}
This repository implementation is a prototype, not a production RTE solver. It
uses periodic pseudo-spectral streaming, uniform directions on
\(\mathbb{S}^1\), normalized isotropic scattering, and a constant-coefficient
periodic diffusion augmentation. The experimental extension below also includes
a constant-anisotropy phase matrix and a Fourier-space reference
transport-scattering TSA prototype. It does not model upwind sweeps or physical
inflow, vacuum, and reflecting boundaries. It also does not model positivity
preservation, spatially varying anisotropy, general phase functions, or
transport-consistent boundary-compatible DSA/TSA. Its purpose is to test the
algebraic NPBS extension and the role of alternative \(B_h\) choices.
\end{remark}

\section{Hard RTE Regimes and Synthetic-Acceleration Experiments}

\subsection{Why the Hard Cases Are Hard}

Write the discrete-ordinates RTE as
\[
  A\psi=(T-S)\psi=q,
  \qquad
  (T\psi)_\ell=\Omega_\ell\cdot\nabla\psi_\ell+\sigma_t\psi_\ell,
  \qquad
  (S\psi)_\ell=\sigma_s\sum_m w_m p_{\ell m}\psi_m .
\]
Classical source iteration is
\[
  \psi^{n+1}=T^{-1}(S\psi^n+q),
  \qquad
  e^{n+1}=T^{-1}Se^n .
\]
The periodic transport-Born prototype replaces \(T^{-1}\) by a constant
reference resolvent \(G_\eta=L_\eta^{-1}\), so
\[
  e^{n+1}=G_\eta V_\eta e^n,\qquad V_\eta=L_\eta-A .
\]
Thus raw transport Born iteration is a Neumann iteration. It becomes slow or
unstable when the scattering and coefficient-mismatch perturbation
\(V_\eta\) is large in the \(G_\eta\)-metric.

In a homogeneous isotropic medium, a low-frequency scalar-flux error has the
source-iteration amplification
\[
  \lambda_{\rm SI}(\xi)
  =
  \sigma_s\sum_\ell
  \frac{w_\ell}{\sigma_t+\ii\Omega_\ell\cdot\xi}.
\]
For small \(|\xi|\),
\[
  1-\lambda_{\rm SI}(\xi)
  =
  \frac{\sigma_a+D|\xi|^2}{\sigma_t}
  +O\!\left(\frac{|\xi|^4}{\sigma_t^4}\right),
  \qquad
  \sigma_a=\sigma_t-\sigma_s,\qquad
  D\simeq\frac{1}{d\sigma_t}.
\]
In optically thick high-albedo scaling, \(\sigma_t=O(\varepsilon^{-1})\) and
\(\sigma_a=O(\varepsilon)\), so \(1-\lambda_{\rm SI}=O(\varepsilon^2)\).
The iteration count grows like \(O(\varepsilon^{-2})\). The slow error is a
smooth, nearly isotropic scalar-flux mode.

For anisotropic scattering, the phase matrix has angular eigenvalues
\(\beta_j\). At zero spatial frequency the angular-mode amplification is
approximately
\[
  \lambda_j(0)=c\,\beta_j,\qquad c=\sigma_s/\sigma_t .
\]
Forward-peaked scattering makes several \(\beta_j\) close to one. Then the
slow subspace includes current and higher angular moments, not only scalar
flux. Scalar DSA can remove the \(j=0\) diffusion mode while leaving these
angular modes poorly corrected.

For heterogeneous and mixed-regime media, the constant-reference Born
perturbation contains
\[
  (V_\eta\psi)_\ell
  =
  (\sigma_{\rm ref}+\eta-\sigma_t(x))\psi_\ell
  +S\psi_\ell .
\]
A crude local contraction indicator is
\[
  \frac{
    |\sigma_{\rm ref}+\eta-\sigma_t(x)|+\sigma_s(x)
  }{
    \sigma_{\rm ref}+\eta
  } .
\]
Sharp contrast, near-void channels, high local albedo, and thin/thick
interfaces can drive this ratio close to or above one. Increasing \(\eta\) can
damp the iteration, but it also changes the residual geometry and over-smooths
the correction that must be learned or accelerated.

\subsection{Experimental Setup}

The hard-case experiments use the same periodic pseudo-spectral RTE prototype
as the verifier. The angular phase matrix is the normalized circular
Poisson-kernel model
\[
  p_{\ell m}(g)
  \propto
  \frac{1-g^2}{1-2g\cos(\theta_\ell-\theta_m)+g^2},
  \qquad
  \sum_m w_m p_{\ell m}(g)=1 .
\]
The test grid is \(32\times32\), with \(24\) directions and \(dx=1/32\).
The source is either an isotropic Gaussian
\[
  q_\ell(x,y)=\exp[-100((x-0.32)^2+(y-0.47)^2)]
\]
or the same Gaussian multiplied by a beamed angular factor proportional to
\(\exp(2\cos\theta_\ell)\). The methods are:
\[
\begin{array}{ll}
\text{Born} & B_h=G_\eta,\\[1mm]
\text{periodic scalar DSA} &
  B_h=G_\eta+\mathcal I(-D_0\Delta+\sigma_{a,0})^{-1}
  \Phi(I-AG_\eta),\\[1mm]
\text{RHS-projected reference TSA-H2} &
  B_h=G_\eta+A_{\rm ref}^{-1}\Pi_{|j|\le2}(I-AG_\eta),\\[1mm]
\text{full reference transport-scattering SA} &
  B_h=G_\eta+A_{\rm ref}^{-1}(I-AG_\eta).
\end{array}
\]
Here \(A_{\rm ref}\) is the constant-coefficient reference
transport-scattering operator
\[
  (A_{\rm ref}\psi)_\ell
  =
  \Omega_\ell\cdot\nabla\psi_\ell
  +(\sigma_{\rm ref}+\eta)\psi_\ell
  -\sigma_{s,\rm ref}\sum_m w_m p_{\ell m}(g)\psi_m ,
\]
inverted in Fourier space.  The projection \(\Pi_{|j|\le2}\) is applied to the
post-transport residual before the full reference inverse is used.  Therefore
TSA-H2 is an RHS-projected reference diagnostic, not a fully discretized
reduced angular transport system \(P_{\rm lo}A_{\rm lo}^{-1}R_{\rm lo}\);
streaming can regenerate higher angular modes after the projected residual is
inverted.  The full reference transport-scattering SA is also not a variational
upper bound or production TSA method. It is a diagnostic for the benefit of a
more transport-like synthetic correction. The scalar DSA baseline is
intentionally simple: it uses a constant isotropic diffusion coefficient and
does not use the anisotropic transport cross section
\(\sigma_{\rm tr}=\sigma_t-\sigma_s\beta_1\).

For stationary iteration,
\[
  \psi^{m+1}=\psi^m+\omega B_h(q-A\psi^m),
\]
the relaxation \(\omega\) is selected from
\(\{0.2,0.4,0.6,0.8,1.0\}\), with a limit of \(120\) iterations and tolerance
\(10^{-6}\).  The Krylov experiment uses right-preconditioned flexible GMRES
with the same \(B_h\), restart length \(30\), maximum \(45\) iterations, and
the same tolerance.  The script is
\[
  \texttt{rte/run\_hard\_case\_experiments.py}.
\]

The four hard cases are:
\begin{enumerate}
\item \textbf{Diffusive:}
\(\sigma_t=20\), \(c=0.99\), \(g=0\). This tests the scalar
diffusion-limit mode.
\item \textbf{Heterogeneous:}
\(\sigma_t=1+9b(x,y)+2.5{\bf 1}_{\sin(2\pi y)>0.35}\),
\(c\in[0.82,0.97]\), \(g=0\). This tests constant-reference mismatch.
\item \textbf{Anisotropic:}
\(\sigma_t=5\), \(c=0.96\), \(g=0.85\). This tests forward-peaked angular
modes.
\item \textbf{Mixed:}
\(\sigma_t\approx0.35\) in a thin streaming region and
\(\sigma_t\approx16\) in a thick scattering region, with \(c=0.20/0.985\) and
\(g=0.65\). This tests thin/thick coupling.
\end{enumerate}
The bump \(b(x,y)=\exp[-70((x-0.58)^2+(y-0.50)^2)]\).

\begin{table}[t]
\centering
\scriptsize
\begin{tabular}{llrrrrrr}
\toprule
Case & Method & \(\omega\) & Stat. it. & Stat. final & Stat. fac. & FGMRES it. & FGMRES final\\
\midrule
Diffusive & Born & 1.0 & 120 & \(8.4{\times}10^{-2}\) & 0.980 & 24 & \(7.6{\times}10^{-7}\)\\
Diffusive & DSA & 1.0 & 13 & \(7.0{\times}10^{-7}\) & 0.336 & 8 & \(4.9{\times}10^{-7}\)\\
Diffusive & Ref-H2 & 1.0 & 119 & \(1.0{\times}10^{-6}\) & 0.890 & 10 & \(8.8{\times}10^{-8}\)\\
Diffusive & Ref-full & 1.0 & 119 & \(1.0{\times}10^{-6}\) & 0.890 & 9 & \(8.1{\times}10^{-7}\)\\
Heterogeneous & Born & 1.0 & 120 & \(9.2{\times}10^{-7}\) & 0.891 & 23 & \(6.7{\times}10^{-7}\)\\
Heterogeneous & DSA & 1.0 & 35 & \(8.9{\times}10^{-7}\) & 0.672 & 19 & \(6.5{\times}10^{-7}\)\\
Heterogeneous & Ref-H2 & 1.0 & 42 & \(8.2{\times}10^{-7}\) & 0.716 & 18 & \(9.0{\times}10^{-7}\)\\
Heterogeneous & Ref-full & 1.0 & 40 & \(8.6{\times}10^{-7}\) & 0.705 & 13 & \(3.7{\times}10^{-7}\)\\
Anisotropic & Born & 1.0 & 120 & \(2.0{\times}10^{-3}\) & 0.949 & 23 & \(5.0{\times}10^{-7}\)\\
Anisotropic & DSA & 1.0 & 69 & \(9.9{\times}10^{-7}\) & 0.818 & 20 & \(6.0{\times}10^{-7}\)\\
Anisotropic & Ref-H2 & 1.0 & 33 & \(7.5{\times}10^{-7}\) & 0.652 & 13 & \(4.0{\times}10^{-7}\)\\
Anisotropic & Ref-full & 1.0 & 33 & \(7.5{\times}10^{-7}\) & 0.652 & 8 & \(5.6{\times}10^{-8}\)\\
Mixed & Born & 1.0 & 120 & \(3.3{\times}10^{-3}\) & 0.953 & 45 & \(4.1{\times}10^{-4}\)\\
Mixed & DSA & 1.0 & 120 & \(4.9{\times}10^{-4}\) & 0.938 & 45 & \(3.1{\times}10^{-4}\)\\
Mixed & Ref-H2 & 1.0 & 120 & \(2.7{\times}10^{-4}\) & 0.934 & 45 & \(9.6{\times}10^{-5}\)\\
Mixed & Ref-full & 1.0 & 120 & \(6.3{\times}10^{-6}\) & 0.905 & 45 & \(8.5{\times}10^{-7}\)\\
\bottomrule
\end{tabular}
\caption{Hard-regime RTE benchmark. ``Stat. fac.'' is the measured geometric
factor over the stationary history. FGMRES uses the same method as a right
preconditioner.}
\label{tab:rte-hard-cases}
\end{table}

\subsection{Interpretation of the Results}

The diffusive case shows the classical pathology: raw Born/source iteration
has a factor \(0.98\) and remains at \(8.4\times10^{-2}\) after \(120\)
steps. DSA converges in \(13\) stationary steps and gives the best GMRES
preconditioner. This is exactly the scalar diffusion-limit mode predicted by
the expansion of \(\lambda_{\rm SI}(\xi)\). The reference transport-scattering
correction is not superior here because the missing low-order physics is scalar
diffusion, and the DSA correction is the right asymptotic model.

The heterogeneous case shows moderate reference mismatch. Raw Born eventually
reaches tolerance but slowly. DSA reduces the stationary iteration count from
\(120\) to \(35\), while the full reference correction gives the best FGMRES
iteration count. The
deep reason is that the constant transport Green operator sees the whole
coefficient variation as a perturbation. DSA corrects the smooth scalar part,
but a transport-like correction better captures interface and angular
transmission effects inside Krylov.

The anisotropic case exposes the scalar nature of DSA. Raw Born stalls at
\(2.0\times10^{-3}\). DSA helps, but still needs \(69\) stationary steps.  The
RHS-projected and full reference transport-scattering corrections converge in
\(33\) stationary steps, and the full reference correction reduces FGMRES to
\(8\) iterations. The reason is spectral: with forward-peaked scattering,
angular eigenvalues \(\beta_j\) for \(j\ge1\) remain close to one. Scalar DSA
only attacks the \(j=0\) moment. A transport-scattering synthetic correction
can address angular transport modes left in \(\ker\Phi\).

The mixed case is the hardest. Born, DSA, and the RHS-projected reference
correction all stagnate in stationary
iteration and fail to reach GMRES tolerance within \(45\) Krylov iterations.
The full reference correction reaches \(8.5\times10^{-7}\) in FGMRES, but its
stationary iteration still
has a slow factor \(0.905\). This case combines thin streaming, thick
diffusion, an interface layer, and anisotropy. A scalar diffusion correction
over-smooths the thin region and under-represents angular interface defects.
The RHS-projected angular diagnostic helps, but the full reference correction
is needed to provide a robust Krylov preconditioner under this iteration
budget. This is a budgeted experiment, not a proof that the other methods
cannot converge with more FGMRES iterations or better parameters.

\subsection{Why a Learned Component Can Help}

Learning should not replace the known transport sweep or DSA/TSA correction
blindly. The experiments show that the classical preconditioner removes a
specific analytically understood subspace, while the failures are residual
defects caused by coefficient mismatch, anisotropy, interfaces, and mixed
regime coupling. This suggests a learned \emph{defect corrector}:
\[
  M_\theta r
  =
  B_h r
  +
  C_\theta\!\left(
    B_h(I-AB_h)r,\,
    \sigma_t,\sigma_s,\,
    c=\sigma_s/\sigma_t,\,
    g,\,
    \tau_{\rm loc}
  \right),
\]
where \(B_h\) is DSA or TSA and \(\tau_{\rm loc}\) denotes local optical
thickness features. The network sees the residual left after the classical
synthetic accelerator, not the original full solution map.

The metric-matched training objective is
\[
  \mathcal L_{\rm defect}(\theta)
  =
  \mathbb E_r
  \frac{
    \|B_h(A M_\theta r-r)\|_2
  }{
    \|B_h r\|_2
  }.
\]
This trains exactly the residual propagator used in the solver. If the learned
map makes \(B_hAM_\theta\) closer to the identity on rollout residuals, then
the residual-propagator theorem gives contraction in the \(B_h\)-metric.

There are three concrete learnable components worth prioritizing:
\begin{enumerate}
\item \textbf{Learned DSA closure.} Replace the constant diffusion operator by
a coefficient-conditioned low-order operator
\[
  -\nabla\cdot(D_\theta(x)\nabla\phi)
  +b_\theta(x)\cdot\nabla\phi+\sigma_{a,\theta}(x)\phi .
\]
This targets heterogeneous and interface errors while preserving the DSA
structure.

\item \textbf{Learned TSA restriction and lift.} Learn \(R_\theta\) and
\(P_\theta\), or learn a local angular basis, so that the synthetic transport
space includes exactly the angular moments that remain slow for the local
phase function and optical thickness. This targets the anisotropic case.

\item \textbf{Learned local blending.} Use a small gating model to combine
transport, DSA, and TSA corrections:
\[
  B_{\theta}r
  =
  \alpha_T(x)B_T r
  +\alpha_D(x)B_{\rm DSA}r
  +\alpha_S(x)B_{\rm TSA}r .
\]
The gate should depend on local albedo, optical thickness, anisotropy, and
coefficient gradients. This targets the mixed thin/thick case.
\end{enumerate}

The hard-case results explain why learning is plausible but also why it must
be constrained. DSA already cures the homogeneous diffusion-limit pathology,
so learning should not waste capacity relearning diffusion. TSA already cures
some angular defects, so learning should focus on the missing reduced space,
the low-order coefficients, and the local blending. Training data should
include rollout residuals from raw Born, DSA/TSA stationary iteration, and
FGMRES, because random probes do not reproduce the nonnormal residual
trajectories seen in the mixed case.

\section{Per-Sample Learned RTE Preconditioners}

\subsection{Architecture Design}

The first learned RTE experiment deliberately avoids cross-sample
generalization.  Each hard case has one fixed coefficient field, one fixed
source family, and one trained preconditioner.  This isolates whether learning
can improve the residual propagator at all before asking whether a network can
generalize across media.
It also means the experiment is in-sample: the rollout residuals used for
training include residuals from the same source trajectory used in evaluation.
The result is therefore not a generalization claim.

Two architecture families are natural.

\paragraph{Coefficient-conditioned defect CNN.}
This is closest in spirit to the Helmholtz NPBS neural solvers.  Given a base
synthetic preconditioner \(B_h\), form
\[
  z_0=B_h r,\qquad r_1=r-Az_0,\qquad z_1=B_h r_1 .
\]
A convolutional network then predicts
\[
  M_\theta r
  =
  z_0+C_\theta(z_1,\sigma_t,\sigma_s,c,g,\nabla\sigma_t,\ldots).
\]
For a periodic pseudo-spectral RTE benchmark, \(C_\theta\) can use circular
spatial convolutions, angular-channel mixing, or Fourier layers.  This is the
right architecture for eventual generalization.  It can learn interface
corrections, angular high-moment corrections, and local blending between
transport, DSA, and TSA.  Its drawback is that it has many parameters and
requires a careful residual dataset.

\paragraph{Residual-polynomial defect network.}
For the first per-sample experiment, the repository implements a much smaller
network.  Let \(B_h\) be DSA or reduced TSA.  Define
\[
  h_0=B_h r,\qquad e_0=r-Ah_0,
\]
and recursively
\[
  h_j=B_h e_{j-1},\qquad e_j=e_{j-1}-Ah_j,
  \qquad j=1,\ldots,p.
\]
The learned preconditioner is
\[
  M_\theta r=\sum_{j=0}^{p}\alpha_j h_j ,
\]
where the complex scalars \(\{\alpha_j\}\) are trainable.  This is a tiny
unrolled learned preconditioner: it learns a residual polynomial in the base
error propagator
\[
  E_h=I-AB_h.
\]
If \(\alpha_j=1\) for all \(j\), it is simply \(p+1\) inner Richardson
corrections.  Training chooses a better polynomial for the observed residual
subspace.  This architecture is not meant to be the final generalizing model,
and it is not a neural operator in the usual sense. It is a diagnostic that
tests whether learned residual-polynomial correction can reduce the hard modes
left by the chosen base preconditioner.

\subsection{Training Objective}

For a residual training set \(\mathcal R\), the learned polynomial minimizes
the same metric-matched objective:
\[
  \min_{\alpha_0,\ldots,\alpha_p}
  \frac{1}{|\mathcal R|}
  \sum_{r\in\mathcal R}
  \left(
  \frac{\|B_h(A M_\theta r-r)\|_2}{\|B_h r\|_2}
  \right)^2 .
\]
The residual set contains rollout residuals from the base stationary method
and smooth random operator probes of the form \(r=A v\).  In the experiments
below \(p=3\), so each learned preconditioner uses four applications of the
base \(B_h\).  The script is
\[
  \texttt{rte/train\_learned\_hard\_cases.py}.
\]
Strictly speaking, the polynomial experiment consumes the raw residual \(r\),
not only the coordinate \(z=B_hr\).  Therefore the direct contraction statement
should be read as a condition on the actual residual map
\[
  r\mapsto B_h(r-AM_\theta r),
\]
unless \(M_\theta\) is explicitly reparameterized as a function of \(B_hr\)
alone.

\begin{table}[t]
\centering
\scriptsize
\begin{tabular}{llrrrrrrrr}
\toprule
Case & Model & Inner & Stat. it. & Stat. \(B_h\) & Stat. final & FGMRES it. & FGMRES \(B_h\) & FGMRES final\\
\midrule
Diffusive & Poly-DSA & 4 & 4 & 16 & \(2.0{\times}10^{-8}\) & 3 & 12 & \(1.4{\times}10^{-8}\)\\
Diffusive & Poly-Ref-H2 & 4 & 16 & 64 & \(6.7{\times}10^{-7}\) & 7 & 28 & \(2.6{\times}10^{-7}\)\\
Heterogeneous & Poly-DSA & 4 & 8 & 32 & \(4.3{\times}10^{-7}\) & 6 & 24 & \(7.8{\times}10^{-8}\)\\
Heterogeneous & Poly-Ref-H2 & 4 & 9 & 36 & \(2.2{\times}10^{-7}\) & 7 & 28 & \(2.4{\times}10^{-7}\)\\
Anisotropic & Poly-DSA & 4 & 12 & 48 & \(4.1{\times}10^{-7}\) & 7 & 28 & \(3.7{\times}10^{-7}\)\\
Anisotropic & Poly-Ref-H2 & 4 & 7 & 28 & \(1.8{\times}10^{-7}\) & 5 & 20 & \(4.7{\times}10^{-7}\)\\
Mixed & Poly-DSA & 4 & 89 & 356 & \(9.4{\times}10^{-7}\) & 27 & 108 & \(7.3{\times}10^{-7}\)\\
Mixed & Poly-Ref-H2 & 4 & 120 & 480 & \(2.8{\times}10^{4}\) & 23 & 92 & \(7.8{\times}10^{-7}\)\\
\bottomrule
\end{tabular}
\caption{Per-sample learned residual-polynomial defect preconditioners.  The
``Inner'' column is the number of base \(B_h\) applications per learned
preconditioner call; the \(B_h\)-columns count effective base-preconditioner
applications.}
\label{tab:rte-learned-hard-cases}
\end{table}

\subsection{Learned Experiment Interpretation}

The learned polynomial confirms the main hypothesis.  In the diffusive case,
DSA alone already converges in \(13\) stationary steps; Poly-DSA reduces this
to \(4\) outer steps and reduces FGMRES from \(8\) to \(3\) iterations.  This
does not discover new physics; it learns a residual polynomial for the already
correct DSA mode.

In the heterogeneous case, Poly-DSA reduces the stationary count from \(35\) to
\(8\), and Poly-Ref-H2 reduces it to \(9\).  The learned coefficients are
compensating for constant-reference mismatch and for the fact that one
synthetic correction is not enough near sharp coefficient variation.

In the forward-peaked anisotropic case, Poly-Ref-H2 is best: it reaches
stationary tolerance in \(7\) steps and FGMRES tolerance in \(5\) iterations.
This matches the angular-mode explanation.  DSA can help, but the reference
transport-scattering correction sees
the slow angular modes that scalar diffusion misses.

The mixed case is the stress test.  Base DSA and Ref-H2 failed to reach
stationary tolerance in \(120\) iterations.  Poly-DSA reaches tolerance in
\(89\) steps and makes FGMRES converge in \(27\) iterations.  Poly-Ref-H2 is
effective as a flexible-GMRES preconditioner but unstable as a stationary
iteration, with the residual diverging.  This is an important warning: a
learned preconditioner can be useful inside FGMRES even if the stationary map
\(I-AM_\theta\) is not contractive.  For production use, one should either
add a contraction/stability penalty during training or use the learned map only
inside flexible GMRES.

These results should be read as proof of mechanism, not yet as a fair
wall-time win.  Each Poly-DSA or Poly-Ref-H2 application uses four base
preconditioner applications.  For example, diffusive DSA takes \(13\) base
applications, while diffusive Poly-DSA takes \(4\times4=16\) base applications
as a stationary solver and \(3\times4=12\) base applications inside FGMRES.
The learned polynomial improves the residual propagator and sometimes reduces
effective base calls, but a cost-normalized conclusion requires reporting
base-preconditioner calls, wall time, memory, and stability together.  The
most reliable conclusion is that learning can reshape the residual polynomial
on hard residual subspaces; it is not yet evidence that this tiny learned
model is uniformly faster than optimized classical DSA/TSA.

The experiments use true physical residual norms, but they do not yet compare
against an independently computed reference solution error.  For the linear
systems tested here, residual reduction is the right solver criterion; however,
a complete numerical study should also report scalar-flux error, angular-flux
error, current error, and positivity/undershoot diagnostics against a direct or
high-accuracy Krylov solve.

\section{Audit Findings and Reliability of the Evidence}

An independent audit of the mathematics and code gives the following
conservative interpretation.

\begin{enumerate}
\item \textbf{The algebraic framework is sound.}
Using a fixed linear residual-to-correction map \(B_h\), whether it is a Green
operator, scalar DSA, or a synthetic transport correction, is consistent with
the residual-propagator formulation of NPBS.  The verifier confirms the
transport Born identity and the generic residual-coordinate propagation
identity to roundoff for the operators actually implemented.

\item \textbf{The current RTE code is a periodic Fourier prototype.}
The raw transport baseline is not classical variable-coefficient source
iteration with an upwind sweep.  The DSA baseline is not production DSA with
physical boundary conditions, variable coefficients, or anisotropic transport
cross sections.  The reference transport-scattering correction is not a
production TSA discretization.  Conclusions about production RTE solvers require
upwind \(S_N\), physical boundary conditions, and transport-consistent
DSA/TSA.

\item \textbf{RHS-projected Ref-H2 is not a closed low-order angular solve.}
The implementation projects the post-transport residual to angular modes
\(|j|\le2\), then applies the full angular reference inverse.  Because
\(\Omega\cdot\nabla\) couples angular Fourier content at nonzero spatial
frequency, the correction can regenerate higher angular modes.  A true reduced
TSA comparison should implement
\[
  P_K(P_KA_{\rm ref}P_K)^{-1}P_K r
\]
or a genuine coarse-ordinate/moment transport solve.

\item \textbf{The learned experiment is in-sample.}
Training residuals include rollout residuals from the same coefficient field,
source, and zero-initialized trajectory used in evaluation.  The learned
polynomial results show that per-instance residual-polynomial shaping can help;
they do not show cross-source, cross-coefficient, or cross-grid
generalization.

\item \textbf{Iteration counts are insufficient without cost.}
The learned polynomial uses four base-preconditioner applications per outer
call.  Reference transport-scattering corrections store and apply dense
per-Fourier-mode angular inverses.  A fair solver comparison must report
operator calls, preconditioner calls, setup time, solve time, memory, and
residual-versus-time curves.

\item \textbf{The present metrics are solver metrics, not full transport
diagnostics.}
Relative residuals are useful, but future experiments should also report
quadrature-weighted residuals, scalar-flux error, angular-flux error, current
error, conservation defects, positivity/undershoot, and comparison with an
independent direct or high-accuracy Krylov reference.

\item \textbf{The strongest current conclusion is narrow but promising.}
For the periodic Fourier toy problems, scalar DSA cures the homogeneous
diffusion-limit mode; reference transport-scattering correction helps angular
and mixed-regime defects; and learned residual polynomials can reshape hard
residual subspaces.  This supports the direction of learned
synthetic-acceleration defect correction, but it is not yet a demonstration of
a general learned RTE preconditioner.
\end{enumerate}

\subsection{Next Steps}

The next steps should be ordered by what they de-risk.

\begin{enumerate}
\item \textbf{Add independent validation.}
Assemble the discrete operator matrix on tiny grids and compare
\(\texttt{A\_op}\), \(G_\eta\), DSA, and reference transport-scattering
corrections against explicit matrix products.  Add manufactured periodic
solutions and direct-solve reference errors.

\item \textbf{Make the reduced transport correction mathematically closed.}
Implement \(P_K(P_KA_{\rm ref}P_K)^{-1}P_K\), coarse ordinate solves, or moment
transport solves, and compare them with the current RHS-projected diagnostic.

\item \textbf{Run parameter sweeps.}
Vary grid size, angular order, albedo, anisotropy \(g\), \(\eta\), source
shape, and coefficient contrast.  Report mean, worst-case, and failure cases,
not only one representative run.

\item \textbf{Use held-out learning tests.}
Train on some sources/rollout residuals and evaluate on held-out sources,
held-out residual trajectories, and eventually randomized media.  Compare the
learned polynomial with nonlearned equal-cost baselines:
\(\alpha_j=1\), least-squares polynomial, Chebyshev/minimax polynomial, and
extra inner DSA/reference-correction iterations.

\item \textbf{Move from polynomial fitting to local learned defect models.}
After the equal-cost polynomial baselines are understood, train the
coefficient-conditioned defect CNN.  It should predict the correction to
DSA or the chosen reference transport correction, not replace the entire solver,
and it should be evaluated primarily as
a flexible-GMRES preconditioner until a contraction penalty makes stationary
use reliable.

\item \textbf{Approach production transport gradually.}
Introduce upwind \(S_N\) streaming, vacuum/inflow/reflecting boundary
conditions, transport-consistent DSA, and anisotropic DSA based on
\(\sigma_{\rm tr}\).  Only after those baselines are stable should claims be
made about RTE solver performance beyond the periodic Fourier prototype.
\end{enumerate}

\subsection{Next Network to Train}

The polynomial network proves that learning can improve the residual
propagator, but it only learns global scalar coefficients.  The next useful
network is the coefficient-conditioned defect CNN
\[
  M_\theta r
  =
  z_0+z_1
  +
  s\,C_\theta\!\left(z_0,z_1,\sigma_t,\sigma_s,c,g\right),
  \qquad
  z_0=B_hr,\quad z_1=B_h(r-Az_0),
\]
with periodic spatial convolutions and angular-channel mixing.  A conservative
design is:
\[
\begin{array}{ll}
\text{input channels:} &
  \Re z_0,\Im z_0,\Re z_1,\Im z_1,
  \sigma_t,\sigma_s,c,g,\|\nabla\sigma_t\|,\\
\text{body:} &
  4\text{--}6 circular residual convolution blocks with angular-channel mixing,\\
\text{optional global branch:} &
  low angular Fourier modes and low spatial Fourier modes,\\
\text{output:} &
  \Re \delta\psi,\Im \delta\psi,\\
\text{solver map:} &
  M_\theta r=z_0+z_1+s\delta\psi .
\end{array}
\]
This model should be trained first per sample, then across randomized media.
For stationary use, include a rollout loss
\[
  \|B_h(f-A(u+M_\theta r))\|_2/\|B_h r\|_2
\]
and a stability penalty that discourages growth of consecutive residuals.  For
Krylov-only use, train the metric loss and evaluate inside flexible GMRES.

\subsection{Laptop-Scale Defect CNN Test}

The repository now includes a first small-scale implementation of this network
in
\[
  \texttt{rte/train\_defect\_cnn\_small.py}.
\]
The goal is not to claim a trained general-purpose RTE preconditioner, but to
test whether the coefficient-conditioned defect idea is numerically sane on a
problem small enough for a laptop.

The base map is the scalar DSA-style preconditioner \(B_h=B_{\rm DSA}\).  The
network uses the fair two-application baseline as its initialization:
\[
  z_0=B_hr,\qquad z_1=B_h(r-Az_0),\qquad
  M_\theta r=z_0+z_1+s\,C_\theta(z_0,z_1,\chi),
\]
where \(s=0.05\) in the test below and
\[
  \chi=(\widehat{\sigma_t},\sigma_s/\sigma_t,
  \|\nabla_h\sigma_t\|/\overline{\sigma_t},g).
\]
The real and imaginary parts of \(z_0\) and \(z_1\) are used as residual
features.  The coefficient features gate circular convolution blocks, while the
final projection is initialized to zero.  Hence before training the learned
preconditioner is exactly DSA2, not an unfairly stronger learned map.

The loss is the metric-matched defect objective with a small physical residual
term:
\[
  \frac{\|B_h(AM_\theta r-r)\|_2^2}{\|B_hr\|_2^2}
  +0.2\frac{\|AM_\theta r-r\|_2^2}{\|r\|_2^2}
  +10^{-5}\frac{\|M_\theta r\|_2^2}{\|B_hr\|_2^2}.
\]
Training residuals are five DSA rollout residuals from the source trajectory
plus five smooth random operator probes \(r=Av\).  The laptop run used a
\(20\times20\) periodic grid, \(12\) angles, two residual convolution blocks,
width \(16\), Adam for \(80\) steps with learning rate \(2\times10^{-3}\), and
\(\eta\)-scale \(0.1\).  Evaluation used the training source and one held-out
source with a different position/angular profile.  DSA2 means two inner DSA
defect-correction applications; CNN-DSA2 has the same two \(B_h\) calls plus
the small CNN correction.

\begin{table}[t]
\centering
\scriptsize
\begin{tabular}{lllrrrr}
\toprule
Case & Split & Method & Stat. it. & Stat. \(B_h\) & Stat. final & FGMRES final\\
\midrule
Diffusive & train & DSA & 60 & 60 & \(5.5{\times}10^{-6}\) & \(4.9{\times}10^{-6}\)\\
Diffusive & train & DSA2 & 60 & 120 & \(5.3{\times}10^{-6}\) & \(5.0{\times}10^{-6}\)\\
Diffusive & train & CNN-DSA2 & 60 & 120 & \(4.9{\times}10^{-6}\) & \(5.3{\times}10^{-6}\)\\
Diffusive & heldout & DSA & 60 & 60 & \(5.0{\times}10^{-6}\) & \(4.7{\times}10^{-6}\)\\
Diffusive & heldout & DSA2 & 60 & 120 & \(4.8{\times}10^{-6}\) & \(4.5{\times}10^{-6}\)\\
Diffusive & heldout & CNN-DSA2 & 60 & 120 & \(4.8{\times}10^{-6}\) & \(4.6{\times}10^{-6}\)\\
Anisotropic & train & DSA & 60 & 60 & \(3.4{\times}10^{-5}\) & \(1.9{\times}10^{-5}\)\\
Anisotropic & train & DSA2 & 60 & 120 & \(3.1{\times}10^{-6}\) & \(3.2{\times}10^{-6}\)\\
Anisotropic & train & CNN-DSA2 & 60 & 120 & \(3.0{\times}10^{-6}\) & \(3.0{\times}10^{-6}\)\\
Anisotropic & heldout & DSA & 60 & 60 & \(3.4{\times}10^{-5}\) & \(2.1{\times}10^{-5}\)\\
Anisotropic & heldout & DSA2 & 60 & 120 & \(3.0{\times}10^{-6}\) & \(3.2{\times}10^{-6}\)\\
Anisotropic & heldout & CNN-DSA2 & 60 & 120 & \(3.3{\times}10^{-6}\) & \(3.3{\times}10^{-6}\)\\
Mixed & train & DSA & 60 & 60 & \(5.1{\times}10^{-3}\) & \(5.0{\times}10^{-3}\)\\
Mixed & train & DSA2 & 60 & 120 & \(7.2{\times}10^{-4}\) & \(3.4{\times}10^{-4}\)\\
Mixed & train & CNN-DSA2 & 60 & 120 & \(9.4{\times}10^{-4}\) & \(2.6{\times}10^{-4}\)\\
Mixed & heldout & DSA & 60 & 60 & \(5.2{\times}10^{-3}\) & \(4.7{\times}10^{-3}\)\\
Mixed & heldout & DSA2 & 60 & 120 & \(7.3{\times}10^{-4}\) & \(2.9{\times}10^{-4}\)\\
Mixed & heldout & CNN-DSA2 & 60 & 120 & \(1.0{\times}10^{-3}\) & \(2.3{\times}10^{-4}\)\\
\bottomrule
\end{tabular}
\caption{Laptop-scale coefficient-conditioned defect CNN test.  All FGMRES
runs used \(25\) iterations, so the final residual is the more informative
quantity here.}
\label{tab:defect-cnn-small}
\end{table}

The small CNN is stable under the DSA2 initialization and scaled correction.
The training loss decreased from \(6.5\times10^{-3}\) to \(4.1\times10^{-4}\)
in the diffusive case, from \(1.4\times10^{-1}\) to \(4.7\times10^{-2}\) in the
anisotropic case, and from \(2.1\times10^{-1}\) to \(5.8\times10^{-2}\) in the
mixed case.  The solver impact is modest.  It slightly improves the
training-source DSA2 result for the anisotropic case and improves FGMRES in the
mixed case, including on the held-out source.  It does not uniformly improve
stationary iteration; in the mixed case the stationary residual is worse than
DSA2 even though FGMRES is better.

The conclusion is therefore specific: the coefficient-conditioned defect CNN is
a promising flexible-GMRES preconditioner component, but at this small scale it
is not yet a robust replacement for equal-cost extra DSA steps.  The next
technical improvements should be a contraction penalty for stationary use,
training on multiple held-out residual trajectories, coefficient batches rather
than one medium at a time, and a true reduced TSA/anisotropic-DSA base
preconditioner for the angular hard modes.

\section{Implementation Added}

The repository now contains a minimal RTE generalization in
\(\texttt{rte/models.py}\). It implements a periodic two-dimensional
discrete-ordinates operator and its constant-coefficient transport Green
operator. It also implements a
transport-plus-diffusion synthetic preconditioner, a constant-reference
transport-scattering TSA-style preconditioner, the Born residual identity for
the transport reference, and preconditioned Richardson solve logic. The verifier
\[
  \texttt{rte/verify\_rte\_npbs.py}
\]
checks the identity
\[
  G_\eta(V_\eta\psi+q)-\psi
  =
  G_\eta(q-A\psi)
\]
and verifies residual decay for both the raw transport preconditioner and the
transport-plus-diffusion and transport-plus-TSA preconditioners on a
moderate-scattering case. For synthetic preconditioners it verifies the generic
residual-coordinate propagation identity rather than claiming a Born split.
The hard-regime experiment script is
\[
  \texttt{rte/run\_hard\_case\_experiments.py}.
\]
The per-sample learned-preconditioner experiment script is
\[
  \texttt{rte/train\_learned\_hard\_cases.py}.
\]
The laptop-scale coefficient-conditioned defect CNN experiment script is
\[
  \texttt{rte/train\_defect\_cnn\_small.py}.
\]

\end{document}
